\documentclass[12pt,letterpaper]{article}
\usepackage[T1]{fontenc}
\usepackage{mathptmx}
\usepackage{courier}
\usepackage[letterpaper,margin=1in]{geometry}
\usepackage{amsmath,amssymb}
\usepackage{microtype}
\usepackage{booktabs}
\usepackage{array}
\usepackage{enumitem}
\usepackage{xr}
\externaldocument[N-]{commutation-transitive-groups}
\usepackage[hidelinks]{hyperref}

\newcommand{\F}{\mathbb{F}}
\newcommand{\C}{\mathbb{C}}
\newcommand{\Z}{\mathbb{Z}}
\newcommand{\triv}{\mathbf{1}}
\DeclareMathOperator{\PGL}{PGL}
\DeclareMathOperator{\PSL}{PSL}
\DeclareMathOperator{\GL}{GL}
\DeclareMathOperator{\SL}{SL}
\DeclareMathOperator{\Ind}{Ind}
\DeclareMathOperator{\Cent}{Cent}
\DeclareMathOperator{\Conj}{Conj}

\title{Revision log\\[4pt]\large ``The representation theory of finite
commutation transitive groups'', Draft 3\\(corrected and expanded expository
version)}
\author{David Joyner\\\texttt{wdjoyner@gmail.com}}
\date{November 17, 2026}

\begin{document}
\maketitle

\noindent This log records every change from the original 2000 manuscript
to Draft~1, from Draft~1 to Draft~2, and from Draft~2 to Draft~3. Theorem numbers of the current draft
are generated from the note's \texttt{.aux} file, so compile the note first.
Numbers of the 2000 manuscript and of Drafts~1 and~2 are historical and are
typed by hand.

\section{Files and drafts}\label{sec:files}

\begin{itemize}[itemsep=1pt]
\item All files are in the folder \texttt{commutation-transitive} of the
  repository \url{https://github.com/wdjoyner/alg-graph-thry}.
\item Original: \texttt{comm\_tr6.tex}, ``preliminary version'' dated
  8-5-2000, unpublished. Draft~1 kept a frozen copy in \texttt{archive/}. At
  the author's request, \texttt{archive/} was removed from the repository on
  November 17, 2026. The original remains in the git history, in commit
  \texttt{dd94a1d}, at the path
  \nolinkurl{commutation-transitive/archive/comm_tr6_2000-08-05.tex}. The
  command \texttt{git show dd94a1d:<path>} prints it.
\item Draft 1 (November 16, 2026): commits \texttt{dd94a1d} and
  \texttt{bd9e320}, merged as pull request~\#4.
\item Draft 2 (November 17, 2026): commit \texttt{d05ec8c}. Its starting
  point was the author's hand-edited copy of Draft~1.
\item Draft 3 (November 17, 2026): \texttt{commutation-transitive-groups.tex}
  and its PDF.
\item Referee reports: on the original, \texttt{referee/2026-11-16/}. On the
  author's hand-edited copy of Draft~1, \texttt{referee/2026-11-17/}.
\item Script: \texttt{verification/comm\_tr\_checks.py} (Python~3, standard
  library and NumPy), with its output in
  \texttt{verification/checks\_output.txt}. A full run takes about one
  minute. It is exhaustive for the groups it treats. It was not changed for
  Draft~2, and it was rerun with identical output. Draft~3 added its
  section~13 (the Frobenius group of order~$21$).
\end{itemize}

\section{AI usage}\label{sec:ai}

All AI work was done by Claude (Anthropic, model Claude Opus 5.5). No other
AI system was involved. AI review is not an independent certification of
the results.

\paragraph{Draft 1 (November 16, 2026).}
\begin{enumerate}[itemsep=1pt]
\item \emph{Referee.} Claude refereed \texttt{comm\_tr6.tex} for correctness,
  readability and repairability, and wrote the referee report.
\item \emph{Reviser and co-author of new material.} At the author's request,
  Claude turned the paper into an expository note. It wrote all new text,
  proofs, examples, the literature review, the notation table and the
  glossary. The house style sheet \texttt{joyner-math-paper-style-guide.md}
  (Draft~2) was applied.
\item \emph{Computation.} Claude wrote and ran \texttt{comm\_tr\_checks.py}
  (Section~\ref{sec:computations}).
\item \emph{Literature.} Claude ran web searches to check bibliographic data
  (Section~\ref{sec:sources}).
\item \emph{Graph theory section.} Asked whether the subject connects to
  graph theory, Claude proposed studying the class map on the commuting
  graph as a harmonic morphism. The author asked for it to be added, and
  Claude wrote Section~\ref{N-sec:harmonic} and Subsection~\ref{N-sec:harmbg}
  and extended the script.
\end{enumerate}

\paragraph{Draft 2 (November 17, 2026).}
\begin{enumerate}[itemsep=1pt]
\item \emph{Referee.} Claude refereed the author's hand-edited copy of
  Draft~1 for correctness and readability. The report is in
  \texttt{referee/2026-11-17/}.
\item \emph{Reviser.} Claude made the changes listed in
  Section~\ref{sec:d1d2}(b), moved all comparisons with the 2000 manuscript
  out of the note into Section~\ref{sec:moved} of this log, and updated the
  repository.
\item \emph{Literature.} Claude located the published version of the
  ``commutative transitive kernel'' paper and the page range of
  Fine--Rosenberger 2009 (Section~\ref{sec:sources}).
\end{enumerate}

\paragraph{Draft 3 (November 17, 2026).} At the author's request, Claude
rewrote the glossary, added Example~\ref{N-ex:f21} and
Figures~\ref{N-fig:k4}--\ref{N-fig:octa}, added the repository to the
Acknowledgments, and extended the script (Section~\ref{sec:d2d3}).

\section{Changes from Draft 2 to Draft 3}\label{sec:d2d3}

All changes were requested by the author on November 17, 2026.
\begin{itemize}[itemsep=1pt]
\item \emph{Glossary.} Rewritten so that each entry gives a short,
  self-contained definition. Draft~2 had several entries that only pointed
  to the body (``See Definition~\dots''), which defeats the purpose of a
  glossary. Entries were added for the class map, complete graphs and naive
  conjectures. The glossary no longer contains cross-references.
\item \emph{Figures.} Three TikZ figures illustrate how a harmonic morphism
  collapses a large graph onto a small one. In each figure, a vertex of the
  small graph and its fiber have the same color, and vertical edges are
  dashed.
  Figure~\ref{N-fig:k4} accompanies Example~\ref{N-ex:harmbg} (a harmonic
  and a non-harmonic collapse of $K_4$). Figure~\ref{N-fig:f21} accompanies
  the new Example~\ref{N-ex:f21}. Figure~\ref{N-fig:octa} accompanies
  Example~\ref{N-ex:octahedron} (the octahedron folded onto a triangle).
  Example~\ref{N-ex:harmbg} and Example~\ref{N-ex:octahedron} were expanded
  to describe their figures.
  Example~\ref{N-ex:a5harm} now refers to Figure~\ref{N-fig:k4}(a).
\item \emph{New Example~\ref{N-ex:f21}}: the class map of the Frobenius
  group $C_7\rtimes C_3$ of order~$21$. One component, $K_6$, collapses onto
  an edge with fibers of size~$3$ and degree~$9$. The seven other
  components map isomorphically. Proved in the text, and checked by the new
  section~13 of the script.
\item \emph{Acknowledgments.} Now name the GitHub repository that holds the
  revision log, the script and the referee reports.
\item \emph{Title.} Subtitle changed to ``Draft 3''.
\end{itemize}
Adding Example~\ref{N-ex:f21} shifts by one the numbers of the later items
in Section~\ref{N-sec:harmonic}. For example, Draft~2's
Example~9.6, Proposition~9.7, Corollary~9.8 and Example~9.9 are now
Example~\ref{N-ex:a5harm}, Proposition~\ref{N-prop:noncommharm},
Corollary~\ref{N-cor:noncommct} and Example~\ref{N-ex:octahedron}. No
theorem, proposition or lemma changed its mathematical content.

\section{Changes from Draft 1 to Draft 2}\label{sec:d1d2}

\paragraph{(a) Changes made by the author, in his hand-edited copy.}
\begin{itemize}[itemsep=1pt]
\item Removed Counterexample~6.2 (the 2000 induced-character formula),
  Counterexample~7.3 (the 2000 proof of the main theorem) and
  Remark~2026.1 (the GAP origin of the bound~$4$). Their content is
  preserved in Section~\ref{sec:moved}.
\item Removed several references to the 2000 manuscript from the body, and
  reduced the literature-review paragraph on the original draft to one
  sentence.
\item Reworded the abstract into three paragraphs, and the introduction's
  ``What this note does'' paragraph.
\item Introduced informal ``Conjecture'' statements before
  Counterexamples~\ref{N-cex:pgl}, \ref{N-cex:quotient} and~\ref{N-cex:pglcor},
  before Proposition~\ref{N-prop:psl2odd}, before
  Proposition~\ref{N-prop:taudim}, and at the start of
  Section~\ref{N-sec:span}.
\item Reworded Counterexample~\ref{N-cex:quotient} and the opening of
  Section~\ref{N-sec:perm}.
\end{itemize}

\paragraph{(b) Changes made by Claude.} Each change answers an item of the
Draft~2 referee report (numbers in brackets).
\begin{itemize}[itemsep=1pt]
\item Subtitle changed to ``Draft 2'', and the title footnote dated
  [S1].
\item Repaired three broken references (\verb|cex:formula|,
  \verb|cex:supports|, \verb|rem:gap|) left by the removals, in the
  introduction and in the summary table [C1].
\item Removed the remaining comparisons with the 2000 manuscript from the
  body: the introduction, Example~\ref{N-ex:s3conj},
  Counterexample~\ref{N-cex:pgl}, the summary table and the
  Acknowledgments [S2]. They are recorded in Section~\ref{sec:moved}.
\item Turned the informal conjectures into numbered ``Naive conjecture''
  statements (Naive Conjectures~\ref{N-conj:pgl}, \ref{N-conj:psl},
  \ref{N-conj:quotient}, \ref{N-conj:pgltori}, \ref{N-conj:tau}
  and~\ref{N-conj:span}). Each has a precise statement and a short motivation
  [R1]. The introduction now says which are false and which are true [C3].
\item Made Naive Conjecture~\ref{N-conj:quotient} precise (``if $G$ is
  commutation transitive, then so is $G/Z(G)$'') [C4].
\item Made Naive Conjecture~\ref{N-conj:span} precise. The draft's version,
  about ``linear combinations of monomial representations'', is true for
  every finite group by Brauer's theorem, and used the undefined symbol
  $G^*$ [C5].
\item Corrected the justification in Question~\ref{N-q:profinite} [C2].
\item Moved Wu's quotient-closure statement into Remark~\ref{N-rem:wu},
  since it rests on secondary sources, and rewrote it so that it says
  something non-trivial [C6].
\item Retitled Subsection~\ref{N-sec:malnormal} [C7]. Reworded the sentence
  introducing Propositions~\ref{N-prop:frobct} and~\ref{N-prop:psl2even}
  [R3].
\item Bibliography: the ``commutative transitive kernel'' paper is now
  cited from its publication, \texttt{FGRS97}, and described in the
  pre-2000 literature review. Added the page range of \texttt{FR09}. Removed
  the uncited entry \texttt{Jo00}. Softened the claim about Suzuki's paper
  [C8].
\item Added orienting paragraphs to the subsections that lacked them (style
  sheet, rule 3.2) [S3]. Fixed the typo ``draafted'' by merging that
  sentence into the opening of the literature review [R2].
\end{itemize}
No theorem, proposition or lemma changed its mathematical content.

\section{Material moved out of the note in Draft 2}\label{sec:moved}

The style sheet keeps revision history out of the body (rule~6.4). The
following comparisons with the 2000 manuscript appeared in Draft~1 and are
recorded here instead.

\subsection{The coset map (Example~\ref{N-ex:s3conj})}
The 2000 manuscript stated the bijection between $\Conj(g,G)$ and the cosets
of $\Cent(g,G)$ with \emph{left} cosets, $x^{-1}gx\leftrightarrow
x\Cent(g,G)$. That map is not well defined. Example~\ref{N-ex:s3conj} still
shows why, without the historical remark.

\subsection{\texorpdfstring{$\PGL(2,\F_q)$ (Counterexample~\ref{N-cex:pgl})}{PGL(2,q)}}
The 2000 manuscript (its Example~2) claimed that $\PGL(2,\F_q)$ is CT for odd
$q$. Its own footnote to Remark~2 contained the pair of matrices that refutes
the claim.

\subsection{The induced-character formula (Draft 1, Counterexample 6.2)}
The 2000 manuscript (its Lemma~8) stated the formula of
Proposition~\ref{N-prop:indformula}(c) with the sum over $N_G(C)$ and without
the factor $|C|^{-1}$. It also gave $|N_G(C)|\mu(g)$ for central $g\neq1$,
instead of $[G:C]\mu(g)$. Draft~1 refuted it as follows. Take $G=S_3$,
$C=A_3$ and $\mu$ non-trivial. At a $3$-cycle the true value is $-1$, while
the 2000 formula gives $\omega+\omega^2$ summed over $N_G(C)=S_3$, which is
$-3$.

\subsection{The proof of the main theorem (Draft 1, Counterexample 7.3)}
The 2000 manuscript (its Theorem~11) argued that functions ``supported on
pairwise distinct sets'' cannot satisfy a linear relation. Draft~1 refuted
this principle. Take three classes $\alpha,\beta,\gamma$, with
$f_1=\triv_\alpha+\triv_\beta$, $f_2=\triv_\alpha+\triv_\gamma$ and
$f_3=2\cdot\triv_\alpha+\triv_\beta+\triv_\gamma$. The supports are
distinct, and even independent over $\F_2$. Yet $f_3=f_1+f_2$. The converse
fails too. In $A_5$ the four vectors of Example~\ref{N-ex:a5tau} are
dependent over $\F_2$, but the four permutation characters are linearly
independent over $\C$ (Table~\ref{N-tab:a5perm}).
Proposition~\ref{N-prop:indep} is the corrected statement.

\subsection{The conjectured bound 4 (Draft 1, Remark 2026.1)}
The 2000 manuscript (its Conjecture~12) was based on GAP computations for
$\PSL(2,\F_{2^k})$ with $2\le k\le5$, which suggested the bound~$4$. Draft~1
guessed that those computations counted the distinct vectors ($4$) rather
than the rank ($3$). The original GAP code has not been available to check
this. In Draft~2 the conjecture appears as Naive
Conjecture~\ref{N-conj:tau}, and Proposition~\ref{N-prop:taudim} proves it
with the sharp bound~$3$. Draft~2 reuses the label ``Remark~2026.1'' for a
different statement (Remark~\ref{N-rem:wu}).

\subsection{The closing claim (Section~\ref{N-sec:span})}
The 2000 manuscript ended by asserting that its main theorem shows that all
irreducible representations can be obtained from linear combinations of
monomial representations, and attributed this kind of result to Gelfand and
Graev. That did not follow from the theorem. Theorem~\ref{N-thm:span} is a
correct version. The relevant tool is Brauer's induction theorem, not the
Gelfand--Graev representations.

\subsection{Literature and acknowledgments}
\begin{itemize}[itemsep=1pt]
\item Draft~1 described the original draft in the literature review and
  cited a companion preprint, D.~Joyner, ``Representation theory of $A_5$
  using GAP'' (2000), which has not been located. Draft~2 keeps one sentence
  about the year 2000, because rule~2.3 of the style sheet splits the
  literature review at the date of the original.
\item The 2000 manuscript cited the commutative transitive kernel paper as
  ``to appear in \emph{Algebra Colloquium}''. Draft~2 cites the published
  version.
\item The 2000 manuscript thanked Tony Gaglione, Dennis Spellman and Jill
  McGowan. Draft~2 keeps the thanks without the historical remark.
\end{itemize}

\subsection{The summary table of Draft 1}
Draft~1's summary table had a column ``Relation to the 2000 draft''. Draft~2
replaces it with a column describing the content. The removed column read as
follows.

\begin{center}
\renewcommand{\arraystretch}{1.1}
\begin{tabular}{>{\raggedright}p{5.4cm}>{\raggedright\arraybackslash}p{8.6cm}}
\toprule
Result (Draft 2 number) & Relation to the 2000 manuscript\\
\midrule
Theorem~\ref{N-thm:classif} & Restated in full, attributions corrected\\
Proposition~\ref{N-prop:gl2} & Was cited without proof\\
Counterexample~\ref{N-cex:pgl}, Proposition~\ref{N-prop:psl2odd} & The 2000
Example~2 was false\\
Lemma~\ref{N-lem:cent}, Proposition~\ref{N-prop:selfnorm} & Corrected
versions of 2000 statements\\
Propositions~\ref{N-prop:frobct}, \ref{N-prop:psl2even} & New\\
Proposition~\ref{N-prop:indformula} & The 2000 formula was off by $|C|$\\
Theorem~\ref{N-thm:irred}, Corollary~\ref{N-cor:frobchars} & New\\
Proposition~\ref{N-prop:pglnorm} & Replaces the 2000 Corollary~9\\
Proposition~\ref{N-prop:indep} & Replaces the 2000 main theorem\\
Proposition~\ref{N-prop:taudim} & Settles the 2000 conjecture (dimension
$3$, not $4$)\\
Theorem~\ref{N-thm:span} & Correct form of the 2000 closing claim\\
Proposition~\ref{N-prop:ccc}, Theorem~\ref{N-thm:harmonic} & New (graph
theory)\\
Proposition~\ref{N-prop:noncommharm}, Corollary~\ref{N-cor:noncommct} & New
(graph theory)\\
\bottomrule
\end{tabular}
\end{center}

\section{Original results and new results}\label{sec:split}

\paragraph{From the 2000 manuscript, kept (corrected where noted).}
\begin{itemize}[itemsep=1pt]
\item The definitions of CT and commutation transitive groups.
\item The statement of the classification (Theorem~\ref{N-thm:classif}). It
  was extended to the full classification, and the attributions were
  corrected.
\item Commutation transitivity of $\GL(2,k)$. It was cited in 2000 and is
  now proved (Proposition~\ref{N-prop:gl2}).
\item The commuting graph lemma (Proposition~\ref{N-prop:graph}) and the
  centralizer lemma (Lemma~\ref{N-lem:cent}). Part (c) of
  Lemma~\ref{N-lem:cent} corrects ``disjoint or equal''.
\item The self-normalizing centralizer observation
  (Proposition~\ref{N-prop:selfnorm}), corrected to pass to $G/Z$.
\item The equivalence with abelian centralizers
  (Proposition~\ref{N-prop:basic}).
\item The induced-character formula (Proposition~\ref{N-prop:indformula}),
  corrected by the factor $|C|^{-1}$ and at central elements.
\item The support map $\tau$ (Definition~\ref{N-def:supp}).
\item Weisner's $k>2$ remark (Remark~\ref{N-rem:weisnerk}).
\item The bound in the 2000 conjecture, now Naive
  Conjecture~\ref{N-conj:tau}.
\end{itemize}

\paragraph{New in 2026.}
\begin{itemize}[itemsep=1pt]
\item Proposition~\ref{N-prop:frobct} (Frobenius groups are CT) and
  Proposition~\ref{N-prop:psl2even} (structure of $\PSL(2,\F_{2^f})$).
\item Proposition~\ref{N-prop:psl2odd} ($\PSL(2,\F_q)$, $q$ odd, is CT only
  for $q=3,5$).
\item Theorem~\ref{N-thm:irred} (irreducibility criterion) and
  Corollary~\ref{N-cor:frobchars}.
\item Proposition~\ref{N-prop:pglnorm} (Mackey norms for tori in
  $\PGL(2,\F_q)$).
\item Proposition~\ref{N-prop:indep} (independence of permutation
  characters).
\item Proposition~\ref{N-prop:taudim} (dimension of the $\tau$-span).
\item Lemma~\ref{N-lem:nilp} and Theorem~\ref{N-thm:span} (spanning theorem
  via Brauer induction).
\item The graph theory: Lemma~\ref{N-lem:collapse},
  Proposition~\ref{N-prop:ccc}, Lemma~\ref{N-lem:free},
  Theorem~\ref{N-thm:harmonic}, Proposition~\ref{N-prop:noncommharm} and
  Corollary~\ref{N-cor:noncommct}. The 2000 manuscript defined the commuting
  graph (its Lemma~4) but did not use it further.
\item All counterexamples, examples, naive conjectures, the literature
  review, the notation table, the glossary and
  Questions~\ref{N-q:psl}--\ref{N-q:profinite}.
\end{itemize}

\paragraph{From the 2000 manuscript, withdrawn.}
\begin{itemize}[itemsep=1pt]
\item Example~2 ($\PGL(2,\F_q)$ is CT). False. It is now Naive
  Conjecture~\ref{N-conj:pgl}, refuted by
  Counterexample~\ref{N-cex:pgl}.
\item Corollary~9. False for $q=2$, and true for $q\ge3$ for a different
  reason. It is now Naive Conjecture~\ref{N-conj:pgltori}, settled by
  Counterexample~\ref{N-cex:pglcor}.
\item The proof of Theorem~11. Invalid (Section~\ref{sec:moved}).
\item The closing claim about monomial representations. It did not follow.
  It is replaced by Theorem~\ref{N-thm:span}.
\item The general (locally compact) Mackey criterion, Lemma~7 and Remark~4.
  They are not needed for finite groups and were imprecise. The finite
  criterion is Theorem~\ref{N-thm:mackey}(b).
\item Section~5 (profinite groups). It had no content. It is now
  Question~\ref{N-q:profinite}.
\item The paragraph on the method of little groups (\S4.1).
  Corollary~\ref{N-cor:frobchars} covers the solvable CT case directly.
\end{itemize}

\section{Concordances}\label{sec:concordance}

\subsection{From the 2000 manuscript to Draft 2}

\begin{center}
\renewcommand{\arraystretch}{1.12}
\begin{tabular}{>{\raggedright}p{5.2cm}>{\raggedright\arraybackslash}p{9cm}}
\toprule
2000 manuscript & Draft 2\\
\midrule
Theorem 1 (classification) & Theorem~\ref{N-thm:classif} and
Remark~\ref{N-rem:weisnerk}\\
Remark 1 (Weisner, Wu) & Remark~\ref{N-rem:weisnerk} and
Section~\ref{N-sec:lit}\\
Example 2 ($\PGL(2,\F_q)$) & Naive Conjecture~\ref{N-conj:pgl},
Counterexample~\ref{N-cex:pgl}\\
Theorem 3 (free groups, $\GL_2$) & Example~\ref{N-ex:ctexamples}(5),
Proposition~\ref{N-prop:gl2}\\
Remark 2 (relation on $G/Z$) & Naive Conjecture~\ref{N-conj:quotient},
Counterexample~\ref{N-cex:quotient}\\
Lemma 4 (commuting graph) & Proposition~\ref{N-prop:graph}\\
Lemma 5 (centralizer lemma) & Lemma~\ref{N-lem:cent}\\
Paragraph on $C=N_G(C)$ & Proposition~\ref{N-prop:selfnorm},
Example~\ref{N-ex:selfnorm}\\
Lemma 6, Remark 3 (abelian centralizers) & Proposition~\ref{N-prop:basic}\\
Lemma 7, Remark 4 (Mackey criterion) & Theorem~\ref{N-thm:mackey}\\
Lemma 8 (Frobenius formula) & Proposition~\ref{N-prop:indformula}, and
Section~\ref{sec:moved} of this log\\
Corollary 9 ($\PGL_2$ tori) & Proposition~\ref{N-prop:pglnorm}, Naive
Conjecture~\ref{N-conj:pgltori}, Counterexample~\ref{N-cex:pglcor}\\
Definition 10 (support) & Definition~\ref{N-def:supp}\\
Theorem 11 (main theorem) & Proposition~\ref{N-prop:indep}, and
Section~\ref{sec:moved} of this log\\
Conjecture 12 & Naive Conjecture~\ref{N-conj:tau},
Proposition~\ref{N-prop:taudim}\\
Closing claim (monomial reps) & Naive Conjecture~\ref{N-conj:span},
Theorem~\ref{N-thm:span}, Example~\ref{N-ex:a5span},
Counterexample~\ref{N-cex:q8}\\
Section 5 (profinite) & Question~\ref{N-q:profinite}\\
\bottomrule
\end{tabular}
\end{center}

\subsection{From Draft 1 to Draft 2}

Sections~1--4 and~9--11 keep their numbering. The changes are in
Sections~5--8.

\begin{center}
\renewcommand{\arraystretch}{1.12}
\begin{tabular}{>{\raggedright}p{6.2cm}>{\raggedright\arraybackslash}p{8cm}}
\toprule
Draft 1 & Draft 2\\
\midrule
Subsection 5.3, ``Centralizers are malnormal'' & Subsection~\ref{N-sec:malnormal},
retitled\\
Subsection 5.5, ``Three false claims of the 2000 draft'' & Subsection~\ref{N-sec:false}, retitled\\
--- & Naive Conjectures~\ref{N-conj:pgl}, \ref{N-conj:psl},
\ref{N-conj:quotient} (new)\\
Counterexample 5.16 & Counterexample~\ref{N-cex:pgl}\\
Proposition 5.17 & Proposition~\ref{N-prop:psl2odd}\\
Counterexample 5.18 & Counterexample~\ref{N-cex:quotient}\\
--- & Remark~\ref{N-rem:wu} (new; text from Counterexample 5.18)\\
Counterexample 6.2 & Removed (Section~\ref{sec:moved})\\
Example 6.3 & Example~\ref{N-ex:a5ind}\\
Theorem 6.4, Corollary 6.5 & Theorem~\ref{N-thm:irred},
Corollary~\ref{N-cor:frobchars}\\
Example 6.6 & Example~\ref{N-ex:frobchars}\\
Proposition 6.7 & Proposition~\ref{N-prop:pglnorm}\\
--- & Naive Conjecture~\ref{N-conj:pgltori} (new)\\
Counterexample 6.8 & Counterexample~\ref{N-cex:pglcor}\\
Counterexample 7.3 & Removed (Section~\ref{sec:moved})\\
Proposition 7.4 & Proposition~\ref{N-prop:indep}\\
--- & Naive Conjecture~\ref{N-conj:tau} (new)\\
Proposition 7.5 & Proposition~\ref{N-prop:taudim}\\
Remark 2026.1 (GAP bound) & Removed (Section~\ref{sec:moved})\\
--- & Naive Conjecture~\ref{N-conj:span} (new)\\
Theorem 8.1, Lemma 8.2, Theorem 8.3 & Theorem~\ref{N-thm:brauer},
Lemma~\ref{N-lem:nilp}, Theorem~\ref{N-thm:span}\\
Example 8.4, Counterexample 8.5 & Example~\ref{N-ex:a5span},
Counterexample~\ref{N-cex:q8}\\
\bottomrule
\end{tabular}
\end{center}

\section{Sources}\label{sec:sources}

Verification levels follow the style sheet. Level~1 is an independent
derivation. Level~2 is a check against a source actually consulted. Level~3
is reliance on memory.

\paragraph{Consulted (level 2, bibliographic data and abstracts or first
pages only).}
\begin{itemize}[itemsep=1pt]
\item Weisner 1925: title, volume and pages from the AMS website.
\item Suzuki 1957: title (``\emph{odd} order''), volume and pages from the
  AMS website, and the first page from an online copy. The first page states
  condition (W) and that $LF(2,2^n)$ are the only simple groups satisfying
  it. The 2000 manuscript's bibliography said ``even order''. That was an
  error.
\item Brauer--Suzuki--Wall 1958: bibliographic data from the reference list
  of Suzuki 1957 on the AMS website.
\item Schmidt 1970: bibliographic data from Numdam. The description of its
  later use comes from the abstract of a 2010 paper by Herzog and coauthors,
  and from Grazian--Monetta 2023.
\item Wu 1998: bibliographic data and the zbMATH review (via the MaRDI
  portal). The review states the solvable result and the correction of
  Weisner. The statements about locally finite CT groups, Tarski monsters and
  closure under quotients come from secondary summaries (Wikipedia and a
  paper summary site), not from the paper. This is why
  Remark~\ref{N-rem:wu} is marked unverified.
\item Fine--Gaglione--Rosenberger--Spellman 1997 (Draft 2): journal, volume,
  year and pages from G.~Rosenberger's publication list (Universit\"at
  Hamburg) and from the reference list of a paper by Moravec and coauthors.
  B.~Fine's own publication list gives ``No.~2 (1997) 141--152'', which is
  the issue number. One Springer reference list gives volume~2, which
  conflicts with the other two sources. The note follows the majority,
  volume~4. The description of the kernel in Section~\ref{N-sec:lit} follows
  a ResearchGate abstract that quotes the definition. The paper itself was
  not consulted.
\item Fine 1983: bibliographic data from the reference list of the 2015
  arXiv paper by Fine, Gaglione, Rosenberger and Spellman.
\item The 2015 paper by Fine, Gaglione, Rosenberger and Spellman
  (\texttt{arXiv:1506.02636}): introduction and reference list.
\item Fine--Rosenberger 2009: the page range 112--130 (Draft 2) from
  G.~Rosenberger's and B.~Fine's publication lists.
\item Abdollahi--Akbari--Maimani 2006: bibliographic data and abstract.
\item Grazian--Monetta 2023: arXiv abstract and Section 2 excerpt.
\item Duncan--Kazachkov--Remeslennikov 2005: arXiv abstract. The abstract
  states the characterization by centralizer dimension~$2$. The author names
  are from memory (level~3).
\item Dol\v{z}an--Klep--Moravec: preprint first page. Publication details
  were not found.
\item Baker--Norine 2009: journal, issue and pages from the reference list
  of a 2025 arXiv paper (\texttt{2504.08539}). Also the first pages of the
  arXiv version of Baker--Norine itself (\texttt{0707.1309}).
\item Herzog--Longobardi--Maj 2009: authors, title, year and DOI from a
  paper summary site and the introductions of two later arXiv papers. The
  volume and pages were not obtained.
\item Rezaei--Foruzanfar 2023: authors, title, journal, DOI and the scope
  ($G/Z$ Frobenius of order $pq$ or $p^2q$) from a paper summary site and
  the zbMATH review on MaRDI.
\item Corry 2010: arXiv abstract and introduction.
\item Joyner--Melles 2017: bibliographic data from library and bookseller
  records.
\end{itemize}

\paragraph{Not consulted (level 3).}
\begin{itemize}[itemsep=1pt]
\item The full texts of Weisner, Suzuki, Brauer--Suzuki--Wall, Schmidt and
  Wu. Theorem~\ref{N-thm:classif} is cited on the strength of the abstracts,
  reviews and secondary summaries listed above. Its statement agrees across
  all of them.
\item Serre 1977: the section numbers for the induced character formula
  (\S7.2), Mackey (\S7.3--7.4) and Brauer's theorem (Chapter~10) are from
  memory.
\item Isaacs 1976, Theorem~7.2 (Frobenius' theorem): the theorem number is
  from memory.
\item Rotman 1995, Chapter~8 (simplicity of $\PSL(2,\F_q)$): from memory.
\item Suzuki, \emph{Group Theory II} (1986): the page range 291--305 for the
  classification is reported by Wikipedia. It is not cited by page in the
  note.
\item Levin--Rosenberger 1986: not located. The exact title (``commutation''
  or ``commutative'') is unchecked.
\item Baker--Norine 2009, Section~2: the statements of the degree lemma and
  the Riemann--Hurwitz formula (Theorem~\ref{N-thm:bn}) are from memory.
  The note never depends on them. Every degree and every Riemann--Hurwitz
  identity it uses is verified directly in Lemma~\ref{N-lem:collapse}.
\item Urakawa 2000: the bibliographic data are from memory.
\item Melles--Joyner 2026: the data are as recorded by the author.
\item The claim in Section~\ref{N-sec:lit} that Suzuki's paper was an early
  step toward the odd order theorem is standard, and from memory.
\end{itemize}

\section{Computational verification}\label{sec:computations}

\texttt{comm\_tr\_checks.py} checks the following by exhaustive computation
with explicit matrices and exact group tables. The output was identical when
rerun for Draft~2. Draft~3 added item~12.
\begin{enumerate}[itemsep=1pt]
\item Which of $\GL(2,q)$, $\PGL(2,q)$, $\PSL(2,q)$ ($q=3,5,7$) and
  $\SL(2,q)$ ($q=4,8$) have abelian centralizers of non-central elements.
  (Proposition~\ref{N-prop:gl2}, Counterexample~\ref{N-cex:pgl},
  Propositions~\ref{N-prop:psl2even} and~\ref{N-prop:psl2odd}.)
\item The induced-character value in $S_3$ that refutes the 2000 formula
  (Section~\ref{sec:moved}), and the irreducibility at $q=2$
  (Counterexample~\ref{N-cex:pglcor}).
\item The Mackey norms for split tori in $\PGL(2,q)$, $q=3,5,7$, for every
  character, and for the non-split torus in $\PGL(2,3)\cong S_4$.
  (Proposition~\ref{N-prop:pglnorm}.)
\item The $\F_2$-rank of the $\tau$-vectors for $\PSL(2,q)$, $q=4,8,16$.
  (Proposition~\ref{N-prop:taudim}.)
\item The left and right coset maps of Lemma~\ref{N-lem:conjcent} in $S_3$.
  (Example~\ref{N-ex:s3conj}.)
\item The $\C$-rank of the permutation characters of $A_5$, with their
  values. (Table~\ref{N-tab:a5perm}.)
\item The four $\Z$-combinations of Example~\ref{N-ex:a5span} have norm~$1$
  and degrees $3,3,5,4$.
\item Theorem~\ref{N-thm:irred} for every centralizer and every character
  of $S_3$, $A_4$ and $A_5$.
\item Theorem~\ref{N-thm:harmonic}, vertex by vertex, for $S_3$, $A_4$,
  $A_5$, $\PSL(2,\F_8)$, $\GL(2,\F_3)$, $Q_8$ and $D_8$: harmonicity, the
  identities $m(x)=|W(\Cent(x,G))|$ and $v(x)=|W|-1$, and the
  Riemann--Hurwitz identity on each component. The component and image
  sizes agree with Table~\ref{N-tab:harmonic}.
\item Failure of harmonicity of the class map on the non-commuting graph of
  $A_5$, at all $59$ vertices (Example~\ref{N-ex:octahedron}).
\item Proposition~\ref{N-prop:noncommharm} for the same seven groups. The
  map is harmonic exactly for $Q_8$ and $D_8$, where Riemann--Hurwitz gives
  $12=12$.
\item (Draft 3.) Example~\ref{N-ex:f21}: the Frobenius group of order~$21$
  is CT, with class sizes $1,3,3,7,7$. Its class map is harmonic, with
  components $K_6\to K_2$ ($|W|=3$, degree~$9$) and $K_2\to K_2$ ($|W|=1$),
  and the Riemann--Hurwitz identity holds.
\end{enumerate}
No result of the note depends on the computations. All of them are proved by
hand. The computations are independent checks.

\section{Style-sheet notes}\label{sec:style}

\begin{itemize}[itemsep=1pt]
\item The document class \texttt{article} was kept from the original
  (style sheet, default for revisions).
\item The packages \texttt{newtxtext} and \texttt{newtxmath} are not
  installed in the build environment. The fallback \texttt{mathptmx} was
  used, as the style sheet allows.
\item The author line lists only David Joyner. AI contributions are credited
  in the Acknowledgments (style sheet, default).
\item The title page still identifies the 2000 manuscript, in the date line
  and in a \verb|\thanks| footnote. Rules~6.2 and~6.3 of the style sheet
  require this. All other comparisons with the 2000 manuscript were moved
  here (rule~6.4).
\item Citation keys use author initials plus a two-digit year.
\item Semicolons appear only inside the optional arguments of cited
  theorems, as in ``[Mackey; cited from \dots]''.
\item Theorem labels are unchanged from Draft~1, except for the removed
  items and the new naive conjectures. All references are generated by
  \verb|\ref|.
\end{itemize}

\end{document}
